\documentclass[10pt,a4paper]{amsart}
\usepackage[margin=19mm]{geometry}
\usepackage{amsmath,amssymb,mathtools}
\usepackage[scr=rsfs,cal=euler]{mathalfa}
\usepackage{xfrac}
\usepackage[unicode,hidelinks,bookmarks=false]{hyperref}
\newcommand{\ie}{i.e.\ }
\newcommand{\cf}{cf.\ }
\newcommand{\resp}{respectively\ }
\newcommand{\Subsection}{\subsection}
\providecommand{\nobreakdash}{\nobreak\hbox{-}}
\setlength{\emergencystretch}{2em}
\multlinegap0pt
\usepackage{bm}
\usepackage{bbm}
\usepackage{dsfont}
\usepackage{xspace}
\usepackage{cancel}
\usepackage{mathtools}
\usepackage{amsthm}
\makeatletter
\@ifundefined{proposition}{\newtheorem{proposition}{Proposition}}{}
\makeatother
\title[Optimization under rare events: scaling laws for linear chan]{Optimization under rare events: scaling laws for linear chance-constrained programs}
\author{Jose Blanchet, Joost Jorritsma, Bert Zwart}
\date{Source: arXiv:2407.11825v2 (CC BY 4.0)}
\begin{document}
\maketitle

\noindent\emph{Source and licence.} Optimization under rare events: scaling laws for linear chance-constrained programs. Version arXiv:2407.11825v2 (CC BY 4.0), distributed under the Creative Commons Attribution 4.0 International licence, as recorded in the arXiv licence metadata for this exact version and shown on the abstract page https://arxiv.org/abs/2407.11825v2. Full rights notice: licenses/arxiv-2407.11825-CC-BY-4.0.txt.\par\smallskip

\noindent\textit{Editorial scope.} Counted unit: the Proposition whose source label is \texttt{prop:1} in the article (source0.tex lines 714--716, source label \texttt{prop:1}), delivered together with the article's own complete original proof of it (source0.tex lines 717--734, one \texttt{proof} environment of 18 lines). No mathematics is written, repaired or re-derived: statement and proof are the article's own text line for line. Required context retained uncounted: none. Every definition, display and label of those lines is retained unchanged. Omitted from this package and not redistributed: 1--713, 735--1110. Nothing inside a retained excerpt is omitted or altered: every retained body is delivered exactly as the article's lines are written, so no omission receipt of any kind accompanies this package. Citations: the retained excerpts carry no citation command at all, so no bibliography entry of the article is transcribed and no reference is redistributed. Reference adaptation: none was needed and none was made; every reference inside the delivered unit resolves inside it, so no unresolved reference or citation is produced. Printed slips kept literal: no printed word, symbol, hypothesis or number of the counted statement or of its proof was altered. Layout and class: the article's own class is \texttt{amsart} with packages including a4wide, amsmath, amssymb, amsfonts, dsfont, comment, mathtools, amsaddr, enumitem, graphicx, subcaption, geometry, xcolor, setspace; the wrapper is the lane's pinned \texttt{amsart} editorial wrapper at 10pt on A4, which restates the theorem-like environments the retained text needs and adds the article's own macro definitions that the retained text needs (none), with extra wrapper packages (bm, bbm, dsfont, xspace, cancel, mathtools). The delivered numbering is the unit's own. Assets: no retained span contains a figure environment, a graphics-inclusion command, a PDF-page inclusion, a table or a TikZ picture; no image file is copied into this package, the package declares no asset dependency and no image file is redistributed with it. Licence: version arXiv:2407.11825v2 of this article is distributed under the Creative Commons Attribution 4.0 International licence, as recorded in the arXiv licence metadata for this exact version and shown on the abstract page https://arxiv.org/abs/2407.11825v2. A full rights notice is delivered at licenses/arxiv-2407.11825-CC-BY-4.0.txt. Characters: every retained excerpt is pure ASCII, so the delivered engine, XeLaTeX with its default Latin Modern text font, reports no missing character.
\begin{proposition}\label{prop:1}
    Let $g_\ell, g: [0,\infty)^m\mapsto \mathbb{R}$ be functions, such that $y\mapsto g_\ell( y)$ is non-decreasing for each $\ell$, and $\eta\mapsto g(\eta y)$ is strictly increasing for each $y$ as a function of the scalar $\eta\in [0,\infty)$. Let $c\in[0,\infty)^m$, and define $v_\ell:=\sup \{c^Ty: g_\ell(y)\le 1, y\ge 0\}$ and $v_\ast := \sup \{c^Ty: g(y)\le 1, y\ge 0\}$. Then, $v_\ell\to v_\ast$.
\end{proposition}

\begin{proof}
    Assume first that $v_\ast=\infty$. For each fixed $a>0$, there exists $y^{(a)}\ge0$ such that $c^Ty^{(a)} \geq a$ and $g(y^{(a)})\leq 1$.
    Fix $\eta\in(0,1)$. By the pointwise convergence $g_\ell(\eta y^{(a)})\to g(\eta y^{(a)})<1$, where the strict inequality follows from the fact that $g(\eta y^{(a)})$ is strictly increasing in $\eta$. In other words, $\eta y^{(a)}$ is feasible for $g_\ell$ for all sufficiently large $\ell$. Moreover, for such $\ell, v_\ell\ge c^T \eta y^{(a)}=\eta c^Ty^{(a)}\ge \eta a$. Since $a$ is arbitrary, we obtain that $\liminf_{\ell\to\infty} v_\ell =\infty=v_\ast$.
    Assume now that $v_\ast<\infty$. Analogously to above, $v_\ell\ge \eta v_\ast$ for each $\eta\in(0,1)$ and all $\ell$ sufficiently large depending on $\eta$. As a result, $\liminf_{\ell \to\infty} v_\ell \ge v_\ast$. 
    We  turn to the upper bound. 

    Let us partition $c=(\bar{c},0)$, where the part $\bar{c}$ has strictly positive elements. We then note that $v_* = \sup\{\bar{c}^T \bar{y} : g(\bar{y},0) \leq 1, \bar{y} \geq 0\}$. Because $g_\ell$ is  non-decreasing 
    we have that 
$v_\ell = \sup\{\bar{c}^T \bar{y}: g_\ell(\bar{y},0) \leq 1, \bar{y}\geq 0\}.$ 
 So, we may assume without loss of generality that $c$ has strictly positive elements.
 
    Assume for contradiction that there exists $\varepsilon>0$ such that $(v_\ell)_{\ell\ge 1}$ contains a subsequence such that $v_{\ell_i} \ge v_\ast+3\varepsilon$ for all $i$. Since $v_{\ell_i}=\sup \{c^Ty: g_{\ell_i}(y)\le 1, y\ge 0\}$,  there must exist $y_{\ell_i}\in [0,\infty)^m$ such that $g_{\ell_i}(y_{\ell_i})\le 1$ and $c^T y_{\ell_i}\ge v_\ast+2\varepsilon$. For each $i$, there exists $\eta_i\leq 1$ and $y'_{\ell_i}\in[0,\infty)^m$ such that $y_{\ell_i}'=\eta_i y_{\ell_i},  g(y'_{\ell_i}) = g(\eta_i y_{\ell_i}) \le g(y_{\ell_i}) \le 1$ and $c^Ty'_{\ell_i}=v_\ast+2\varepsilon$. 


    Let $A:=\{y\in[0,\infty)^n: c^Ty=v_\ast + 2\varepsilon\}$. As all elements of $c$ are strictly positive, $A$ is a compact set. 
    Define $f(y):=(\lfloor y_1\cdot |c|/\varepsilon\rfloor, \ldots, \lfloor y_n\cdot |c|/\varepsilon\rfloor)\cdot \varepsilon/|c|$, i.e., each coordinate is rounded down to the nearest multiple of $\varepsilon/|c|$. Define $\hat A:=\{f(y): y\in A\}\ni f(y'_{\ell_i})=: \hat y_{\ell_i}$. Since $A$ is compact, $\hat A$ is a finite set, and thus we can extract a further subsequence such that $\hat y_{\ell_{i_j}}$ is identical to some $\hat y\in \hat A$ along the subsequence. By the rounding procedure, $c^T \hat y\ge v_\ast+\varepsilon$, and since $g_{\ell_{i_j}}$ is non-decreasing, $g_{\ell_{i_j}}(\hat y)\le 1$. 
        By the pointwise convergence, $g(\hat y)=\lim_{j\to\infty}g_{\ell_{i_j}}(\hat y)\le 1$, and thus $\hat y$ is a feasible solution for $g$. This contradicts that $v_\ast$ is the supremum of $\{c^Ty: g(y)\le 1, y\ge 0\}$. As a result, $\limsup_{\ell\to\infty}v_\ell\le v_\ast+3\varepsilon$ for any $\varepsilon>0$. Thus, $\limsup_{\ell\to\infty}v_\ell\le v_\ast$. 
\end{proof} 

\end{document}
